\section{Aplicaciones del aprendizaje por refuerzo en el mundo real}

Esta sección presenta tres aplicaciones importantes del aprendizaje por refuerzo profundo: el RL de extremo a extremo en la conducción autónoma, AlphaGo y AlphaZero en el go, y, en años recientes, el RLHF en la alineación de los LLM.

\subsection{Conducción autónoma: de la simulación al despliegue real}

\begin{figure}[H]
\centering
\includegraphics[width=0.95\textwidth]{slides-images/slide-48.jpg}
\caption{VISTA: simulador de conducción autónoma de alta fidelidad desarrollado por el MIT\protect\footnotemark}
\end{figure}
\footnotetext{Diapositiva 48. Amini+ IEEE RA-L 2020.}

El profesor Amini presentó su trabajo de investigación en el MIT, en el que usa el método de gradiente de política para entrenar sistemas de conducción autónoma. Las innovaciones clave de este trabajo incluyen:

\textbf{El simulador VISTA}: un entorno de simulación de alta fidelidad basado en datos reales de conducción. A diferencia de la simulación tradicional por gráficos computacionales, VISTA usa video real de conducción para sintetizar nuevos puntos de vista, lo que garantiza que las imágenes simuladas tengan un realismo fotográfico.

\textbf{Entrenamiento de RL de extremo a extremo}:
\begin{itemize}
    \item El estado del agente es la imagen de la cámara de conducción
    \item La acción es el ángulo continuo del volante
    \item El diseño de la recompensa es extremadamente disperso: solo se otorga una penalización cuando el vehículo se sale del carril o sufre una colisión
    \item La red de política produce los parámetros de una distribución gaussiana, de la cual se muestrea el ángulo de dirección concreto
\end{itemize}

\begin{figure}[H]
\centering
\includegraphics[width=0.95\textwidth]{slides-images/slide-49.jpg}
\caption{El primer automóvil autónomo del mundo entrenado completamente mediante RL en simulación y desplegado directamente en carreteras reales\protect\footnotemark}
\end{figure}
\footnotetext{Diapositiva 49. Amini+ IEEE RA-L 2020.}

\begin{knowledgebox}{De la simulación a la realidad (Sim-to-Real Transfer)}
Lo innovador de este trabajo es que el modelo se entrena \textbf{por completo en un entorno simulado}, sin haber tenido nunca contacto con datos reales de carretera, pero puede \textbf{desplegarse directamente en un automóvil real} y conducir con seguridad. Este es el primer automóvil autónomo de tamaño real del mundo entrenado por completo mediante RL en simulación y desplegado en el mundo real. Lo esencial está en que la alta fidelidad del simulador VISTA cierra la brecha entre la simulación y la realidad (Domain Gap).
\end{knowledgebox}

\begin{figure}[H]
\centering
\includegraphics[width=0.95\textwidth]{slides-images/slide-47.jpg}
\caption{La dificultad de aplicar el entrenamiento de RL en el mundo real: ejecutar la política (paso 2) puede provocar peligros\protect\footnotemark}
\end{figure}
\footnotetext{Diapositiva 47.}

\begin{warningbox}{El riesgo de entrenar RL en entornos reales}
El paso 2 del algoritmo de entrenamiento por gradiente de política, "ejecutar la política hasta la terminación", puede ser sumamente peligroso en un entorno real. En el caso de la conducción autónoma, "terminación" puede significar una colisión o un accidente. Por ello, el \textbf{entorno simulado} es el requisito previo para que el RL se aplique en dominios críticos para la seguridad: el agente puede cometer errores y aprender con seguridad en la simulación, sin causar daño real.
\end{warningbox}

\subsection{El go: AlphaGo y AlphaZero}

El go (Go) es una de las aplicaciones más llamativas del RL profundo. La complejidad del go supera con mucho a la del ajedrez: el número de posiciones legales en un tablero de 19x19 es de aproximadamente $2.08 \times 10^{170}$, más que la cantidad de átomos del universo.

\begin{figure}[H]
\centering
\includegraphics[width=0.95\textwidth]{slides-images/slide-52.jpg}
\caption{El espacio de búsqueda del go: el número de posiciones legales supera la cantidad de átomos del universo\protect\footnotemark}
\end{figure}
\footnotetext{Diapositiva 52. Fuente: Wikipedia.}

\textbf{AlphaGo} (Silver+ Science 2016) empleó un proceso de entrenamiento en varias etapas:

\begin{figure}[H]
\centering
\includegraphics[width=0.95\textwidth]{slides-images/slide-54.jpg}
\caption{El proceso de entrenamiento de AlphaGo: de los datos de expertos humanos al autojuego\protect\footnotemark}
\end{figure}
\footnotetext{Diapositiva 54. Silver+ Science 2016.}

\begin{enumerate}
    \item \textbf{Preentrenamiento supervisado}: usar partidas de expertos humanos para entrenar la red de política (Policy Network) como una tarea de clasificación (predecir la jugada de un jugador humano)
    \item \textbf{Autojuego con aprendizaje por refuerzo}: usar el método de gradiente de política del RL para que la red de política siga mejorando mediante el autojuego, hasta alcanzar un nivel \textbf{superior al humano}
    \item \textbf{Entrenamiento de la red de valor}: usar los datos generados por el autojuego para entrenar la red de valor (Value Network), que evalúa la probabilidad de victoria de una posición (tarea de regresión)
\end{enumerate}

En 2016, AlphaGo venció 4 a 1 al campeón mundial de go Lee Sedol, un hito en la historia de la IA.

\textbf{AlphaZero} (Silver+ Science 2018) fue un paso más allá: \textbf{abandonó por completo los datos humanos} y se entrenó desde cero únicamente mediante aprendizaje por refuerzo con autojuego:

\begin{figure}[H]
\centering
\includegraphics[width=0.85\textwidth]{slides-images/slide-57.jpg}
\caption{AlphaZero: aprendizaje desde cero por completo mediante autojuego, con la puntuación Elo subiendo rápidamente conforme avanzan los pasos de entrenamiento\protect\footnotemark}
\end{figure}
\footnotetext{Diapositiva 57. Silver+ Science 2018.}

\begin{importantbox}{El salto de AlphaGo a AlphaZero}
\begin{itemize}
    \item \textbf{AlphaGo}: necesita partidas de expertos humanos como punto de partida (preentrenamiento supervisado + ajuste fino con RL)
    \item \textbf{AlphaZero}: comienza por completo desde cero (tabula rasa), sin usar ningún dato humano, y alcanza, solo mediante \textbf{autojuego más aprendizaje por refuerzo}, un nivel superior al de todas las versiones anteriores
    \item AlphaZero demostró que, en juegos de información perfecta, el autojuego con RL puro puede \textbf{descubrir estrategias que los humanos nunca imaginaron}, superando los límites de la cognición humana
\end{itemize}
\end{importantbox}

\subsection{Resumen de la sección}

El RL profundo ha logrado resultados innovadores en ámbitos como la conducción autónoma y el go. El RL de extremo a extremo en la conducción autónoma muestra la capacidad del gradiente de política para abordar tareas de control continuo, y el simulador es la clave para garantizar un entrenamiento seguro. La serie AlphaGo y AlphaZero demuestra que el RL profundo combinado con el autojuego puede alcanzar, e incluso superar, el nivel humano en problemas de decisión extremadamente complejos.

